\documentclass[conference]{IEEEtran}
\IEEEoverridecommandlockouts

\usepackage{cite}
\usepackage{amsmath,amssymb}
\usepackage{booktabs}
\usepackage{graphicx}
\usepackage{xcolor}
\usepackage{url}

\begin{document}

\title{Decoupling Storage I/O and Memory Management for\
High-Throughput Analytical Scans}

% TODO: replace the anonymous block with the final author list after review.
\author{\IEEEauthorblockN{Anonymous Authors}
\IEEEauthorblockA{Paper under preparation for ICDE}}

\maketitle

\begin{abstract}
Modern NVMe arrays expose enough parallelism that an analytical scan can be
limited by software overhead, buffer ownership, and interference between I/O
completion and decoding rather than by raw device bandwidth.  We study this
problem in Pixels, a columnar storage engine integrated with DuckDB.  Its
current scan path supports synchronous reads, \texttt{io\_uring}, and an SPDK
prototype, but backend selection, asynchronous completion, and buffer lifetime
are spread across the format and physical layers.  We present a format-neutral
I/O and memory-management architecture that separates read planning, backend
sessions, reusable buffer leases, and file-level prefetch pipelines.  The same
protocol accommodates blocking \texttt{pread}, interrupt-driven
\texttt{io\_uring}, and polling-based SPDK without exposing file descriptors,
registered-buffer indices, queue pairs, or logical block addresses to a format
decoder.  An initial implementation includes reusable dynamic and static
buffer pools, single- and double-buffered scans, a Parquet adapter, and
instrumentation that separates I/O waiting from decoding.  Preliminary results
on a 24-NVMe server show that double buffering reduces query time by
6.0--10.0\% for \texttt{io\_uring} and 9.7--17.9\% for SPDK across 12--48
workers.  At 48 workers, however, it increases cache misses by 7.1\% while
instructions remain nearly unchanged, demonstrating that hiding I/O can shift
the bottleneck to memory-intensive decoding.  These observations motivate
backend-aware resource control and asynchronous metadata pipelines rather than
unbounded scan concurrency.
\end{abstract}

\begin{IEEEkeywords}
analytical databases, asynchronous I/O, buffer management, NVMe, io\_uring,
SPDK, columnar storage
\end{IEEEkeywords}

\section{Introduction}

NVMe devices and multi-SSD arrays have changed the balance of analytical query
execution.  A scan operator must keep many devices busy while decoding and
materializing columns, yet every additional worker also creates I/O queues,
buffers, completion work, and pressure on shared caches and memory bandwidth.
Optimizing either storage or computation in isolation can therefore move, and
occasionally worsen, the end-to-end bottleneck.

Pixels is a columnar storage system whose C++ reader is exposed as a DuckDB
table function.  The reader supports direct synchronous I/O, asynchronous
\texttt{io\_uring}, file-level double buffering, and an experimental SPDK
backend.  This breadth makes it a useful case study, but it also reveals a
common systems problem: backend details have leaked into the format reader.
Backend selection occurs in several layers; a buffer identifier may double as
an \texttt{io\_uring} registration index; and metadata still follows a
synchronous chain even when column data is prefetched asynchronously.

This paper asks three questions.  First, what is the smallest interface that
can express the different ownership and completion semantics of \texttt{pread},
\texttt{io\_uring}, and SPDK?  Second, how should buffers be reused without
letting a backend overwrite memory still referenced by a decoder?  Third, when
I/O and decoding overlap, which resource becomes the next scalability limit?

We make the following contributions:

\begin{itemize}
  \item We identify four sources of coupling in an end-to-end columnar scan:
  backend selection, asynchronous capability, buffer registration, and buffer
  lifetime.
  \item We design a format-neutral I/O manager organized around backend
  sessions, explicit read batches, owned buffer views, and a bounded prefetch
  pipeline.
  \item We implement the major mechanisms in the Pixels C++ and DuckDB paths,
  including dynamic and static buffer reuse, double buffering, CPU affinity,
  profiling, an \texttt{io\_uring}-based Parquet path, and an SPDK prototype.
  \item We present preliminary evidence that eliminating visible completion
  waits does not eliminate resource contention: at high concurrency it exposes
  cache- and memory-bound decoding as the dominant cost.
\end{itemize}

This draft deliberately separates measured observations from proposed work.
All numbers in Section~\ref{sec:evaluation} are preliminary and must be rerun
from a single build before submission.

\section{Background and Motivation}

\subsection{The Pixels Scan Path}

A DuckDB worker obtains a file, constructs a Pixels reader, reads metadata,
plans column ranges, issues physical reads, and decodes the returned bytes into
column vectors.  The current physical path is approximately:

\begin{verbatim}
PixelsRecordReader -> RequestBatch -> Scheduler
 -> PhysicalLocalReader -> RandomAccessFile backend
\end{verbatim}

Column data can be submitted in batches.  With double buffering, a worker
decodes the current file while the next file is in flight.  Metadata is
different: locating column chunks requires a strict sequence,

\begin{equation}
\textit{TailOffset} \rightarrow \textit{FileTail}
\rightarrow \textit{RowGroupFooter} \rightarrow \textit{ColumnChunks}.
\end{equation}

The first three stages are currently synchronous.  Consequently, asynchronous
column reads do not remove file-boundary stalls, and SPDK workers can busy-poll
for small metadata requests.

\subsection{Why One Interface Is Difficult}

The backends differ in more than their submission call.  \texttt{pread} uses a
file offset and completes inline.  \texttt{io\_uring} uses submission and
completion queues and may register buffers.  SPDK uses DMA-capable memory,
per-thread queue pairs, polling, and a file-to-LBA mapping.  A useful common
interface must preserve these differences behind a session while providing the
same lifetime and error semantics above it.

\subsection{The Overlap Paradox}

Double buffering should hide I/O latency, but it also allows more workers to
decode simultaneously.  In our preliminary 48-worker run, visible
\texttt{io\_uring} completion waiting fell by 98.3\%, yet accumulated decoding
time rose by 8.0\%.  Instructions increased by only 0.15\%, whereas cycles and
cache misses increased by 4.7\% and 7.1\%, respectively.  Thus, the system did
not decode more data; the same work became more expensive under contention.

\section{Design}

\subsection{Layered Architecture}

Our target architecture has four layers.  A \emph{format adapter} converts
metadata into read ranges.  An \emph{I/O manager} selects a backend and creates
a worker-scoped session.  A \emph{buffer provider} allocates reusable memory
that satisfies alignment, registration, and DMA requirements.  Finally, a
\emph{prefetch pipeline} controls the number and lifetime of in-flight slots.
Pixels and Parquet keep independent decoders but share these lower layers.

\begin{figure}[t]
\centering
\fbox{\parbox{0.92\columnwidth}{\centering
Pixels Adapter \quad Parquet Adapter\\
$\downarrow$\\
I/O Manager $\rightarrow$ Prefetch Pipeline\\
$\downarrow$\hspace{1.3cm}$\downarrow$\\
Buffer Provider \quad I/O Session\\
\hfill Pread $\mid$ io\_uring $\mid$ SPDK}}
\caption{Format-neutral I/O and memory-management architecture.}
\label{fig:architecture}
\end{figure}

\subsection{Owned Buffers and Views}

A \emph{buffer allocation} is the sole owner of a memory region and carries
capacity, alignment, and DMA properties.  A \emph{buffer lease} represents its
temporary use by a pipeline slot.  A read result returns an immutable view that
retains shared ownership of the allocation.  Therefore, a pool cannot recycle
or overwrite the region while an I/O batch or decoder still references it.
Registration tokens and SPDK DMA metadata remain private to the session.

This ownership rule replaces implicit relationships among raw pointers,
buffer-pool entries, and backend request objects.  It also supports different
providers: ordinary aligned allocations, dynamically sized reusable buffers,
pre-sized static rings, and SPDK huge-page allocations.

\subsection{Read Batches and Sessions}

A read range contains a request identifier, file offset, length, target
allocation, and target offset.  A session submits one or more ranges and
returns a batch with the following state machine:

\begin{verbatim}
Created -> Submitted -> Completing -> Completed
                    \-> Failed
Submitted -> Cancelling -> Cancelled | Completed
\end{verbatim}

After submission, the session retains every target until all commands have
completed or have been drained after failure.  Completion order is associated
through request identifiers rather than queue order.  Short reads are errors
unless explicitly permitted.  A blocking backend may return a batch already in
the completed state, preserving the caller protocol without pretending that
all implementations are asynchronous.

\subsection{Bounded Prefetch}

The prefetch pipeline maintains current and next slots.  Each slot owns its
file handle, batch, leases, and decoder-visible views.  The next slot may be
submitted while the current slot is decoded, but it cannot be promoted until
completion.  The current slot cannot be reclaimed until its last decoder view
is released.  The number and total capacity of slots are bounded to prevent
prefetch from turning device parallelism into uncontrolled memory pressure.

\subsection{Asynchronous Metadata}

We propose a metadata coordinator that advances a bounded window of files
through three asynchronous stages: tail-offset reads, file-tail reads, and
row-group-footer reads.  Ready readers enter a worker queue.  A query-level
single-flight cache prevents concurrent misses from issuing duplicate reads.
For SPDK, metadata uses DMA buffers and the same completion infrastructure as
column data; for \texttt{io\_uring}, requests are batched into SQEs.  This is a
design under implementation, not part of the preliminary performance claims.

\section{Implementation Status}

The prototype is implemented in C++ and integrated with DuckDB.  Table
\ref{tab:status} distinguishes completed mechanisms from work needed for a
submission-quality artifact.

\begin{table}[t]
\caption{Prototype status at the time of this draft.}
\label{tab:status}
\centering
\begin{tabular}{p{0.53\columnwidth}p{0.32\columnwidth}}
\toprule
Component & Status \\
\midrule
Pixels vectorized reader/writer and DuckDB scan & Implemented \\
Direct synchronous and \texttt{io\_uring} I/O & Implemented \\
Dynamic/global static buffer pools & Implemented \\
Single/double-buffer scan modes & Implemented \\
Column-vector reuse and CPU affinity & Implemented \\
Parquet \texttt{io\_uring} adapter and profiling & Prototype \\
SPDK raw-block backend and DMA pool & Prototype \\
Unified public manager/session API & Design; partial mechanisms \\
Cross-file asynchronous metadata & Planned \\
Shared SPDK completion pollers & Planned \\
\bottomrule
\end{tabular}
\end{table}

The \texttt{io\_uring} implementation includes fixed and non-fixed buffer
variants.  The global static pool is initialized from collected column-size
information; the dynamic pool adapts allocation sizes at runtime.  The Parquet
adapter preloads column regions into Pixels-managed buffers and exposes them
through an Arrow random-access-file adapter.  Profilers report accumulated
thread time for I/O submission, completion, metadata, decoding, conversion,
and file transitions.  Accumulated thread time is never reported as wall time.

The SPDK path maps immutable Pixels files to raw device extents and reads into
DMA-capable memory.  This prototype currently requires offline mapping and
environment setup; portability, fragmented files, and recovery behavior need
additional engineering before it can be presented as a general backend.

\section{Preliminary Evaluation}
\label{sec:evaluation}

\subsection{Methodology}

The current dataset uses ClickBench query Q24 on a server with 24 Samsung
PM9A3 NVMe SSDs and 12, 24, or 48 scan workers.  We compare single and double
buffering under \texttt{io\_uring} and SPDK.  Wall time, accumulated stage
time, and Linux \texttt{perf} counters are collected.  The SPDK and
\texttt{io\_uring} observations were obtained on different dates; they are
useful for identifying trends but not for declaring a winner between backends.

\begin{table}[t]
\caption{Preliminary Q24 wall time in seconds.}
\label{tab:wall}
\centering
\begin{tabular}{rrrrr}
\toprule
Workers & \multicolumn{2}{c}{io\_uring} & \multicolumn{2}{c}{SPDK} \\
& Single & Double & Single & Double \\
\midrule
12 & 148.90 & 133.97 & 145.67 & 119.55 \\
24 & 77.55  & 70.33  & 74.95  & 64.01  \\
48 & 48.81  & 45.89  & 50.79  & 45.87  \\
\bottomrule
\end{tabular}
\end{table}

\subsection{Double Buffering Hides Completion Wait}

Table~\ref{tab:wall} shows that double buffering improves wall time at all
tested worker counts.  The gain decreases as concurrency grows: 10.0\%, 9.3\%,
and 6.0\% for \texttt{io\_uring}; and 17.9\%, 14.6\%, and 9.7\% for SPDK.  In
the 48-worker \texttt{io\_uring} case, accumulated CQE waiting falls from
266.06 to 4.64 thread-seconds.  In the corresponding historical SPDK runs,
asynchronous completion polling falls from roughly 336 to 7.55 thread-seconds.

\subsection{Decoding Becomes the Bottleneck}

At 48 workers, \texttt{io\_uring} double buffering reduces context switches by
92.6\%, but increases task-clock by 6.8\% and decoding time by 8.0\%.
Instructions are stable, while cache misses rise from 14.006 to 15.002 billion.
This supports a resource-shift interpretation: removing blocked workers
increases the number of simultaneous decoders, enlarges the active working set,
and raises cache and memory-system contention.

\subsection{Synchronous Metadata Remains Visible}

At 48 workers, synchronous metadata accounts for 43.91 accumulated
thread-seconds in the \texttt{io\_uring} single-buffer run and 141.36 in the
SPDK single-buffer run.  About 96 seconds of the latter is synchronous SPDK
polling.  These values are sums across workers and cannot be subtracted from
query wall time.  They nevertheless identify metadata completion as a source
of wasted CPU and motivate the staged pipeline in Section~III-E.

\section{Experimental Plan}

A defensible ICDE evaluation requires a fresh, controlled matrix.  We will use
one build and dataset for all backends, report at least three repetitions with
median and dispersion, and normalize CPU cost by bytes and rows.  The matrix
will include \texttt{pread}, \texttt{io\_uring}, and SPDK; single and double
buffers; 1--24 SSDs; 1--48 workers; Pixels and Parquet; and both ClickBench and
TPC-H queries.  Ablations will isolate buffer reuse, fixed registration,
prefetch depth, CPU affinity, asynchronous metadata, and shared completion
pollers.  Besides wall time and throughput, we will report task-clock, cycles,
instructions, cache misses, context switches, I/O bytes, request latency,
empty-poll ratio, peak memory, and NUMA traffic.

Correctness tests will cover short reads, corrupt metadata, cancellation,
partial batch failure, non-aligned ranges, fragmented files, and the absence of
SPDK hardware.  The primary research comparison is not simply ``which backend
is faster,'' but which combination achieves the best throughput per CPU core
and per byte of reserved memory under each workload regime.

\section{Related Work}

DuckDB demonstrates a vectorized, in-process analytical engine and provides
the execution framework used by our prototype~\cite{duckdb}.  Apache Arrow and
Parquet define widely used in-memory and storage representations; our Parquet
adapter deliberately leaves decoding to Arrow while unifying only raw I/O
buffers~\cite{arrow,parquet}.  \texttt{io\_uring} reduces Linux asynchronous
I/O overhead through shared submission and completion rings~\cite{iouring}.
SPDK moves NVMe drivers into user space and uses polling and DMA memory to
remove kernel-stack overhead~\cite{spdk}.  Our work is complementary: it
focuses on the ownership and scheduling contract required to place these
backends beneath columnar formats, and on the interference created when I/O
latency is successfully hidden.

% TODO: add the published Pixels paper and a broader survey of database buffer
% management, asynchronous prefetching, and computational-storage work.

\section{Conclusion}

Fast storage does not remove I/O management from analytical systems; it makes
the interaction among queues, buffers, completion, and decoding more important.
This draft presents a format-neutral architecture and an initial implementation
in Pixels and DuckDB.  Preliminary experiments show meaningful gains from
overlap, but also show that aggressive overlap can convert I/O stalls into
cache and memory contention.  Completing the common session API, asynchronous
metadata pipeline, and controlled cross-format evaluation will turn the
current engineering prototype into a submission-ready ICDE paper.

\section*{Acknowledgment}

% TODO: funding and acknowledgments.
Omitted for anonymous review.

\bibliographystyle{IEEEtran}
\bibliography{references}

\end{document}
